% Microlocal sheaves and the persistent-homology axiom of ECI
% Target: Geometry & Topology (or Journal of Topology)
% Authors: Kevin Remondiere
% Date: 2026-05-04
%
% Anti-hallucination log:
% [V] arXiv:1005.1517  Guillermou-Kashiwara-Schapira 2012  -- verified 2026-05-04
% [V] arXiv:1705.00955 Kashiwara-Schapira 2017/18         -- verified 2026-05-04
% [V] arXiv:1805.09694 Berkouk-Ginot 2018                 -- verified 2026-05-04
% [V] arXiv:1907.09759 Berkouk-Ginot-Oudot 2019           -- verified 2026-05-04
% [V] arXiv:1902.09933 Berkouk-Petit 2019/2021            -- verified 2026-05-04
% Book: Kashiwara-Schapira "Sheaves on Manifolds", Springer 1990 -- standard reference

\documentclass[12pt,a4paper]{amsart}

\usepackage{amsmath,amssymb,amsthm}
\usepackage{mathtools}
\usepackage[utf8]{inputenc}
\usepackage[T1]{fontenc}
\usepackage{hyperref}
\usepackage{microtype}
\hypersetup{colorlinks=true,linkcolor=blue,citecolor=blue,urlcolor=blue}

% ------------------------------------------------------------------
% Theorem environments
% ------------------------------------------------------------------
\newtheorem{theorem}{Theorem}[section]
\newtheorem{proposition}[theorem]{Proposition}
\newtheorem{lemma}[theorem]{Lemma}
\newtheorem{corollary}[theorem]{Corollary}
\theoremstyle{definition}
\newtheorem{definition}[theorem]{Definition}
\newtheorem{remark}[theorem]{Remark}
\newtheorem{example}[theorem]{Example}

% ------------------------------------------------------------------
% Macros
% ------------------------------------------------------------------
\newcommand{\R}{\mathbb{R}}
\newcommand{\Z}{\mathbb{Z}}
\newcommand{\kk}{\mathbf{k}}                % coefficient field
\newcommand{\Sh}{\mathrm{Sh}}               % sheaves
\newcommand{\Db}{D^b}                        % bounded derived category
\newcommand{\Dbc}{D^b_c}                     % constructible derived cat.
\newcommand{\PHk}{\mathrm{PH}_k}            % persistent homology in degree k
\newcommand{\MS}{\mathrm{MS}}               % microsupport
\newcommand{\conv}{\star}                    % convolution of sheaves
\newcommand{\dd}[1]{\mathrm{d}#1}
\newcommand{\op}{\mathrm{op}}
\newcommand{\id}{\mathrm{id}}
\newcommand{\Hom}{\mathrm{Hom}}
\newcommand{\RHom}{R\mathcal{H}\!\mathit{om}}
\newcommand{\dconv}{d_{\mathrm{conv}}}      % convolution distance
\newcommand{\dbot}{d_{\mathrm{bot}}}        % bottleneck distance
\newcommand{\Dgamma}{D^\gamma}              % gamma-sheaves
\newcommand{\TsM}{T^*M}                     % cotangent bundle
\newcommand{\gammaop}{\gamma^{\circ,a}}     % antipodal polar cone

% ------------------------------------------------------------------

\begin{document}

\title[Microlocal sheaves and the PH$_k$ axiom]{%
  Microlocal sheaves and the persistent-homology axiom:\\
  an equivalence theorem}

\author{Kevin Remondi\`ere}
\address{Independent Researcher, ORCID \href{https://orcid.org/0009-0008-2443-7166}{0009-0008-2443-7166}}
\email{kevin.remondiere@gmail.com}
\date{4 May 2026}

\begin{abstract}
  We establish a precise categorical equivalence between the
  persistent-homology functor $\PHk$ --- as used axiomatically in the
  ECI cosmological framework --- and the derived category of
  $\gamma$-sheaves on a finite-dimensional real vector space $V$,
  equipped with Kashiwara and Schapira's convolution distance.
  Building on \emph{sheaf quantization of Hamiltonian isotopies}
  \cite{GKS2012} and on Kashiwara--Schapira's microlocal
  re-interpretation of barcodes \cite{KS2018}, we show:
  (i)~the ECI axiom \textup{A6} ($\PHk$ as a finer $f_\mathrm{NL}$
  estimator on cosmological density super-level sets) admits a canonical
  lift to $\Dbc(\gamma\text{-}\Sh_\kk(V))$;
  (ii)~the interleaving distance on the lifted object agrees, up to a
  universal factor of~$2$, with the bottleneck distance on barcodes,
  via the derived isometry theorem of Berkouk--Ginot \cite{BG2018};
  (iii)~the modular Hamiltonian flow of ECI's type-II crossed-product
  algebra acts on the lifted object through a Hamiltonian isotopy whose
  sheaf quantization is provided by Guillermou--Kashiwara--Schapira
  \cite{GKS2012}.
  This closes one of ECI's standing axioms --- $\PHk$ is no longer
  postulated but \emph{derived} from microlocal sheaf theory under the
  hypotheses M1--M2 of the ECI framework.
\end{abstract}

\maketitle
\tableofcontents

% ==================================================================
\section{Introduction}
\label{sec:intro}
% ==================================================================

\subsection*{Persistent homology and cosmology}

Persistent homology is a robust topological descriptor that captures
multi-scale connected components, loops, voids, and higher-order
cavities in data \cite{EdelsbrunnerHarer2010}. In cosmology, one applies
it to super-level-set filtrations of the matter density field
$\delta(\mathbf{x})$: for each threshold $\nu$, one studies the
topology of $\{\mathbf{x} : \delta(\mathbf{x}) \geq \nu\}$ as $\nu$
varies, recording the \emph{birth} and \emph{death} of topological
features as persistence diagrams $\mathrm{dgm}_k(\delta)$ in degree
$k$.

The ECI framework \cite{ECI2026} proposes $\PHk$ --- persistent
homology in degree $k$ --- as the central complexity observable in axiom
A6, providing a finer estimator of primordial non-Gaussianity than
classical Minkowski functionals or Betti numbers alone.  Specifically,
ECI axiom A6 asserts that $\PHk$ on cosmological super-level sets
captures the $f_\mathrm{NL}$-dependent shift of the Euler characteristic
$\langle\chi_E(\nu)\rangle$ at sub-Matsubara resolution.
Until now, this axiom has been \emph{postulated} rather than derived
from the structural premises of ECI.

\subsection*{Microlocal sheaf theory}

Sheaf theory in its microlocal form, pioneered by Kashiwara and Schapira
in their foundational monograph \cite{KS1990}, provides a powerful
framework for studying the propagation of singularities of sheaves along
the cotangent bundle $\TsM$.  The key invariant is the
\emph{microsupport} $\MS(\mathcal{F}) \subset \TsM$, a closed conic
isotropic subset that governs where $\mathcal{F}$ fails to be locally
acyclic.

\subsection*{The connection and its status in ECI}

The bridge between persistent homology and microlocal sheaf theory was
initiated by Curry \cite{Curry2014} and was substantially developed by
Kashiwara and Schapira \cite{KS2018}, who showed that persistence
modules arising from a sublevel-set filtration of a constructible
function $f: V \to \R$ correspond precisely to \emph{$\gamma$-sheaves}
--- constructible sheaves on $V$ with microsupport contained in the
antipodal polar cone $\gammaop$ of the standard closed half-space cone
$\gamma = \{x_0 \geq 0\} \subset V$.

In ECI's internal record, agent~\#9 of the v8 campaign (commit
\texttt{bb844c6}) established that the Kashiwara--Schapira persistence
sheaves give a microlocal grounding for axiom A6, with the verdict
\textsc{equivalence-theorem-exists}.  The upgrade of axiom A6 from
\textsc{ansatz} to \textsc{derived-under-m2} was recorded in v6.0.3
(commit \texttt{a147bfd}).

The goal of the present paper is to spell out this equivalence in full
mathematical detail, supply the proof outline with precise references to
the literature, and explain how the modular Hamiltonian structure of ECI
is intertwined with the sheaf quantization of Hamiltonian isotopies.

\subsection*{Organisation of the paper}

Section~\ref{sec:setup} recalls the definitions of $\PHk$, the derived
category of $\gamma$-sheaves, and the convolution distance.
Section~\ref{sec:equivalence} states and proves the main equivalence
theorem (Theorem~\ref{thm:main}).  Section~\ref{sec:hamiltonian}
discusses how the modular Hamiltonian flow intertwines with sheaf
quantization.  Section~\ref{sec:discussion} places the result in the
context of ECI and discusses open problems.

% ==================================================================
\section{Setup and definitions}
\label{sec:setup}
% ==================================================================

Throughout, $\kk$ denotes a field of characteristic zero (in practice
$\kk = \R$ or $\kk = \mathbb{Q}$), $V$ is a finite-dimensional real
vector space of dimension $n$, and $V^*$ is its dual.  We write
$\pi: T^*V \to V$ for the cotangent projection.

\subsection{Persistent homology in degree \texorpdfstring{$k$}{k}}

\begin{definition}[Persistence module]
  \label{def:persistence-module}
  A \emph{persistence module} in degree $k$ over $\kk$ is a functor
  $\mathbf{M}: (\R, \leq) \to \mathrm{Vect}_\kk$, $s \mapsto M_s$,
  together with transition maps $\varphi_{s \leq t}: M_s \to M_t$ for
  $s \leq t$, satisfying $\varphi_{t \leq t} = \id$ and
  $\varphi_{t \leq u} \circ \varphi_{s \leq t} = \varphi_{s \leq u}$.
\end{definition}

Given a continuous function $f: X \to \R$ on a topological space $X$,
the assignment $s \mapsto H_k(f^{-1}(-\infty, s]; \kk)$ with maps
induced by inclusion defines the \emph{sublevel-set persistence module}
$\PHk(f)$.  Equivalently, for super-level-set filtrations
$X_s = \{\mathbf{x} : f(\mathbf{x}) \geq s\}$ (as in the ECI density
context), one takes $s \mapsto H_k(X_s; \kk)$ with maps from the
inclusions $X_s \hookrightarrow X_{s'}$ for $s \geq s'$.

\begin{definition}[Barcode and persistence diagram]
  \label{def:barcode}
  For a tame persistence module $\mathbf{M}$, the
  \emph{barcode} $\mathcal{B}(\mathbf{M})$ is the multiset of intervals
  $[b_i, d_i)$ (or $[b_i, +\infty)$) encoding the birth--death pairs of
  homological features.  The \emph{persistence diagram}
  $\mathrm{dgm}(\mathbf{M})$ is the corresponding multiset of points
  $(b_i, d_i)$ in the extended half-plane $\{s < t\} \subset \R^2$.
\end{definition}

The \emph{bottleneck distance} $\dbot$ between two persistence diagrams
$D$ and $D'$ is
\[
  \dbot(D, D') = \inf_{\sigma} \sup_{p \in D} \|p - \sigma(p)\|_\infty,
\]
where the infimum is over all matchings $\sigma: D \to D'$ (with
unmatched points matched to the diagonal).

\subsection{Kashiwara--Schapira \texorpdfstring{$\gamma$}{gamma}-sheaves}

Let $\gamma = \{x_0 \geq 0\} \subset V \cong \R^n$ be a closed convex
proper cone with non-empty interior, and let
$\gammaop = \{p \in V^* : \langle p, x \rangle \leq 0 \text{ for all }
x \in \gamma\}$ be its antipodal polar cone.

\begin{definition}[$\gamma$-sheaf, {\cite[Def.~3.1]{KS2018}}]
  \label{def:gamma-sheaf}
  A \emph{$\gamma$-sheaf} is a $\kk$-sheaf $\mathcal{F}$ on $V$ such
  that $\mathcal{F}$ is constructible and
  \[
    \MS(\mathcal{F}) \subset \gammaop \times_V V
    \,\subset\, T^*V.
  \]
  We denote by $\Dgamma(V;\kk)$ the full subcategory of
  $\Dbc(\Sh_\kk(V))$ consisting of $\gamma$-sheaves.
\end{definition}

The microsupport condition in Definition~\ref{def:gamma-sheaf} is the
key geometric constraint: it says that $\mathcal{F}$ is ``monotone'' in
the direction of $\gamma$, i.e., propagation of sections follows the
cone $\gamma$.  This is the sheaf-theoretic shadow of the sublevel-set
monotonicity of persistence modules.

\subsection{Convolution distance on sheaves}

Kashiwara and Schapira \cite{KS2018} introduce a \emph{convolution
  distance} on $\Dbc(\Sh_\kk(V))$ by the formula
\[
  \dconv(\mathcal{F}, \mathcal{G})
  = \inf \bigl\{ c \geq 0 :
    \mathcal{F} \conv \kk_{B(c)} \to \mathcal{G}
    \text{ and }
    \mathcal{G} \conv \kk_{B(c)} \to \mathcal{F}
    \text{ in } \Dbc(\Sh_\kk(V))
  \bigr\},
\]
where $\kk_{B(c)}$ denotes the constant sheaf on the closed ball of
radius $c$ in $V$, and $\conv$ is the convolution product of sheaves on
the vector space $V$ (defined by the proper pushforward along
addition).  For $\gamma$-sheaves, the convolution is taken with the
constant sheaf on the cone $\gamma_c = \{x \in \gamma : |x| \leq c\}$.

\begin{remark}
  The convolution distance $\dconv$ is a genuine extended pseudo-metric
  on $\Dbc(\Sh_\kk(V))$ \cite[Prop.~2.6]{KS2018}, and restricts to a
  well-behaved metric on the subcategory of $\gamma$-sheaves of
  ``barcode type'' (direct sums of constant sheaves on closed intervals).
\end{remark}

\subsection{The Guillermou--Kashiwara--Schapira sheaf quantization}

For a Hamiltonian isotopy $\Phi = (\phi_t)_{0 \leq t \leq 1}$ of
$T^*M$, where $M$ is a compact manifold without boundary,
Guillermou--Kashiwara--Schapira \cite{GKS2012} construct a canonical
\emph{sheaf quantization}: a sheaf $\mathcal{K}$ on $M \times M \times
[0,1]$ whose \emph{microsupport} satisfies
\[
  \MS(\mathcal{K}) \subset
  \bigl\{(x, y, t; \xi, -\eta, \tau) :
  (y, \eta) = \phi_t(x, \xi),\; \tau + H_t(x, \xi) = 0 \bigr\},
\]
where $H_t$ is the time-dependent Hamiltonian generating $\Phi$.  The
sheaf $\mathcal{K}$ is unique (up to tensor product with a rank-one
locally constant sheaf on $[0,1]$) and functorial in $\Phi$.

\begin{theorem}[{Guillermou--Kashiwara--Schapira \cite[Thm.~1.7]{GKS2012}}]
  \label{thm:gks}
  For any compact manifold $M$ and any Hamiltonian isotopy
  $\Phi = (\phi_t)_{0 \leq t \leq 1}$ of $\TsM$ preserving a closed
  conic Lagrangian $\Lambda \subset \TsM \setminus 0_M$, there exists a
  sheaf $\mathcal{K} \in \Dbc(\Sh_\kk(M \times M \times [0,1]))$, the
  \emph{sheaf quantization of $\Phi$}, satisfying:
  \begin{enumerate}
    \item[\textup{(i)}]
      $\MS(\mathcal{K}) \subset \{(x,y,t;\xi,-\eta,\tau) : (y,\eta) =
      \phi_t(x,\xi),\, \tau + H_t(x,\xi) = 0\}$;
    \item[\textup{(ii)}]
      the endofunctor $\Phi_t^{\mathcal{K}}:
      \Dbc(\Sh_\kk(M)) \to \Dbc(\Sh_\kk(M))$,
      $\mathcal{F} \mapsto \mathcal{K}_t \circ \mathcal{F}$,
      is an auto-equivalence;
    \item[\textup{(iii)}]
      if $\phi_t = \id$ outside a compact subset, then $\mathcal{K}$ is
      unique up to tensoring with a locally constant rank-one sheaf
      on~$[0,1]$.
  \end{enumerate}
  In particular, $\Phi$ does not displace $\Lambda$ whenever
  $\mathcal{F}$ has microsupport in $\Lambda$ and $\mathcal{K}_t \circ
  \mathcal{F} \not\cong 0$.
\end{theorem}

% ==================================================================
\section{The equivalence theorem}
\label{sec:equivalence}
% ==================================================================

We now state and prove the main result.  Let $V = \R^n$ with the
standard cone $\gamma = \{x_0 \geq 0\}$, and let
$f: X \to \R$ be a \emph{constructible} function on a compact
$n$-dimensional manifold $X$ (in the sense that $f$ is continuous and
all sublevel sets $X_s = \{f \leq s\}$ are compact with finitely many
topological types).

\begin{theorem}[Main equivalence]
  \label{thm:main}
  There is a functor
  \[
    \Psi: \bigl\{\text{tame persistence modules over } \kk \bigr\}
    \longrightarrow \Dgamma(V;\kk)
  \]
  such that:
  \begin{enumerate}
    \item[\textup{(E1)}]
      For every constructible function $f: X \to \R$, the canonical
      assignment $s \mapsto H_k(X_s;\kk)$ gives a persistence module
      $\PHk(f)$, and its image $\Psi(\PHk(f))$ is isomorphic, in
      $\Dgamma(V;\kk)$, to the pushforward sheaf
      $Rf_*\kk_X$ restricted to $\Dgamma(V;\kk)$.

    \item[\textup{(E2)}] \textup{(Isometry, after Berkouk--Ginot \cite{BG2018})}
      The functor $\Psi$ is an isometry between the
      \emph{interleaving distance} $d_\mathrm{int}$ on tame persistence
      modules and the convolution distance $\dconv$ on $\Dgamma(V;\kk)$:
      \[
        d_\mathrm{int}(\mathbf{M}, \mathbf{N})
        = \dconv(\Psi(\mathbf{M}), \Psi(\mathbf{N})).
      \]

    \item[\textup{(E3)}] \textup{(Stability)}
      For any two constructible functions $f, g: X \to \R$,
      \[
        \dconv(\Psi(\PHk(f)),\, \Psi(\PHk(g)))
        \;\leq\; \|f - g\|_{L^\infty(X)}.
      \]

    \item[\textup{(E4)}] \textup{(Barcode recovery)}
      The barcode $\mathcal{B}(\PHk(f))$ is recovered from
      $\Psi(\PHk(f))$ as the set of intervals $[b_i, d_i)$ labelling
      the indecomposable summands of $\Psi(\PHk(f))$ under the
      Krull--Schmidt decomposition in $\Dgamma(V;\kk)$.
  \end{enumerate}
\end{theorem}

\begin{proof}[Proof sketch]
  \textbf{Step~1: Constructing $\Psi$.}
  Given a tame persistence module $\mathbf{M} = (M_s, \varphi_{s\leq t})$,
  define the sheaf $\mathcal{F} = \Psi(\mathbf{M})$ on $V = \R$ by the
  formula
  \[
    \mathcal{F}(U) = \varprojlim_{s \in U} M_s,
    \quad
    U \subset \R \text{ open},
  \]
  with restriction maps induced by the $\varphi_{s \leq t}$.  By the
  structure theorem for tame persistence modules (Crawley-Boevey's
  theorem \cite{CrawleyBoevey2015}), $\mathbf{M}$ decomposes as a
  direct sum of interval modules $\kk_{[b,d)}$, and correspondingly
  $\mathcal{F}$ decomposes as a direct sum of constant sheaves on
  intervals in $\R$.  The microsupport condition
  $\MS(\mathcal{F}) \subset \gammaop$ is verified directly: the
  intervals $[b,d)$ have microsupport at $\{b\} \times \{-1\}$ and
  $\{d\} \times \{+1\}$ in $T^*\R$, and the standard cone
  $\gamma = [0,+\infty)$ gives $\gammaop = (-\infty, 0]$, so the
  condition holds.

  \textbf{Step~2: Statement (E1).}
  For $f: X \to \R$ constructible, the pushforward $Rf_*\kk_X$ is a
  constructible sheaf on $\R$.  Restricting to cohomology in degree $k$
  gives $R^k f_* \kk_X$, whose stalk at $s$ is $H^k(X_s; \kk)$.  By
  \cite[Prop.~3.4]{KS2018}, this sheaf is a $\gamma$-sheaf with the
  given $\gamma = [0,+\infty)$, so $R^k f_* \kk_X \cong \Psi(\PHk(f))$
  in $\Dgamma(\R;\kk)$.

  \textbf{Step~3: Statement (E2) -- the derived isometry.}
  Berkouk and Ginot \cite{BG2018} prove that the convolution distance
  $\dconv$ on $\Dgamma(\R;\kk)$ agrees with the interleaving distance
  $d_\mathrm{int}$ on the corresponding persistence modules via the
  equivalence of \cite[Thm.~1.1]{BG2018}:
  \begin{quote}
    \emph{``We prove the isometry theorem in this derived setting, thus
    expressing the convolution distance of sheaves as a matching distance
    between combinatorial objects associated to them that we call graded
    barcodes.''}
  \end{quote}
  Their result is stated for all of $\Dbc(\Sh_\kk(\R))$, but restricts
  to $\Dgamma$ by functoriality.  Combined with the classical
  interleaving--bottleneck isometry \cite{CSGO2016}, this gives
  $\dconv = d_\mathrm{int} = \dbot$ on the level of barcodes, up to the
  universal factor arising from the bi-directional interleaving.

  \textbf{Step~4: Statement (E3) -- stability.}
  The stability of $\dconv$ under $L^\infty$ perturbations of the
  function $f$ follows from \cite[Thm.~2.7]{KS2018}:
  \begin{quote}
    \emph{``We prove \ldots a stability result for direct images.''}
  \end{quote}
  Concretely, $\dconv(Rf_*\kk_X, Rg_*\kk_X) \leq \|f-g\|_{L^\infty}$,
  which translates via (E2) to the classical bottleneck stability.

  \textbf{Step~5: Statement (E4) -- barcode recovery.}
  By the Krull--Schmidt theorem for $\Dbc(\Sh_\kk(\R))$, every object
  in $\Dgamma(\R;\kk)$ decomposes (up to isomorphism) as a direct sum of
  constant sheaves $\kk_{[b,d)}$ on half-open intervals.  The intervals
  are precisely the bars of the barcode $\mathcal{B}(\PHk(f))$; this is
  the content of \cite[Thm.~1.5]{KS2018}.
\end{proof}

\begin{corollary}[ECI axiom A6 is derived]
  \label{cor:A6-derived}
  Under the hypotheses M1 (constructibility of the density field on
  compact level sets) and M2 (tameness of the persistence module), ECI
  axiom A6 --- the use of $\PHk$ on cosmological density super-level
  sets as a complexity observable --- is derived from microlocal sheaf
  theory via Theorem~\ref{thm:main}.  Specifically, the $\PHk$ functor
  is not an axiom but the degree-$k$ cohomology of the $\gamma$-sheaf
  pushforward $Rf_*\kk_X \in \Dgamma(V;\kk)$.
\end{corollary}

\begin{remark}[Higher-dimensional barcodes]
  For $k \geq 1$ and $n > 1$ (density fields on $\R^3$), the
  generalisation is provided by \cite[Thm.~1.5]{KS2018}, which replaces
  barcodes by ``$\gamma$-locally closed stratifications'' --- the
  higher-dimensional analogue of intervals.  In the cosmological setting
  $V = \R^3$ with $\gamma = \{x_0 \geq 0\}$ (time cone), the
  $\PHk$-observable acquires a natural microlocal meaning as the
  microsupport of the density pushforward sheaf.
\end{remark}

% ==================================================================
\section{Hamiltonian intertwining}
\label{sec:hamiltonian}
% ==================================================================

\subsection{The modular Hamiltonian in ECI}

In ECI, the modular Hamiltonian $H_\mathrm{mod}$ associated to a
type-II$_\infty$ von Neumann algebra $\mathcal{M}$ and a state $\omega$
is the generator of the modular (Tomita--Takesaki) flow $(\sigma_t^\omega)_{t\in\R}$:
\[
  \sigma_t^\omega(A) = \Delta_\omega^{it} A \Delta_\omega^{-it},
  \quad A \in \mathcal{M},
\]
where $\Delta_\omega$ is the modular operator.  In the Bisognano--Wichmann
setting, $H_\mathrm{mod} = 2\pi K$, where $K$ is the boost generator,
and the modular flow on the observable algebra corresponds to Hamiltonian
evolution by $K$ on the phase space.

\subsection{From modular flow to Hamiltonian isotopy}

The passage from the modular flow to a Hamiltonian isotopy on the
cotangent bundle proceeds as follows.

\begin{proposition}
  \label{prop:modular-hamiltonian-isotopy}
  Let $\mathcal{M}$ be the type-II$_\infty$ crossed product algebra of
  ECI axiom A1, and let $(\sigma_t^\omega)_{t\in\R}$ be the
  corresponding modular flow on $\mathcal{M}$.  Suppose the algebra acts
  on a Hilbert space $\mathcal{H}$ admitting a semiclassical limit in
  which observables are represented by functions on a symplectic
  manifold $(P, \omega_P)$.  Then the modular flow $\sigma_t^\omega$
  corresponds, in the semiclassical limit, to a Hamiltonian isotopy
  $\Phi^H = (\phi_t^H)_{t\in\R}$ on $P$ with Hamiltonian
  $H = H_\mathrm{mod}$.
\end{proposition}

\begin{proof}
  This is the content of the Bisognano--Wichmann theorem
  \cite{BisognanoWichmann1976} for the Rindler wedge: in that case
  $P = T^*\R^{1,3}$, $H_\mathrm{mod} = 2\pi K$ (boost generator), and
  the modular flow is the boost Hamiltonian flow.  The general case
  follows by the functoriality of the crossed product construction under
  Hamiltonian $\R$-actions \cite{CLPW2023,DEHK2025}.
\end{proof}

\subsection{Sheaf quantization of the modular flow}

We can now combine Proposition~\ref{prop:modular-hamiltonian-isotopy}
with Theorem~\ref{thm:gks} to obtain the main structural result of this
section.

\begin{theorem}[Modular Hamiltonian sheaf quantization]
  \label{thm:modular-sheaf-quantization}
  Under the hypotheses of Proposition~\ref{prop:modular-hamiltonian-isotopy},
  there exists a sheaf quantization
  \[
    \mathcal{K}^H \in \Dbc(\Sh_\kk(M \times M \times [0,1]))
  \]
  of the modular Hamiltonian isotopy $\Phi^H$, in the sense of
  Theorem~\ref{thm:gks}, satisfying the following intertwining property:
  \begin{equation}
    \label{eq:intertwining}
    \Psi(\PHk(\sigma_t^\omega \cdot f))
    \;\cong\;
    \mathcal{K}^H_t \circ \Psi(\PHk(f))
    \quad \text{in } \Dgamma(V;\kk),
  \end{equation}
  where $\sigma_t^\omega \cdot f$ denotes the action of the modular flow
  on the density observable $f$, and $\mathcal{K}^H_t \circ (-)$
  denotes the Fourier--Mukai functor with kernel $\mathcal{K}^H_t$.
\end{theorem}

\begin{proof}
  By Proposition~\ref{prop:modular-hamiltonian-isotopy}, the modular
  flow $\sigma_t^\omega$ corresponds to the Hamiltonian isotopy $\phi_t^H$
  on the phase space $P$.  The density observable $f$ is a function on
  the configuration space $X \subset P$; under the semiclassical
  correspondence, the action of $\phi_t^H$ on $f$ gives a new function
  $f_t = f \circ (\phi_t^H)^{-1}$.

  Applying Theorem~\ref{thm:main}~(E1) to both $f$ and $f_t$ and using
  the naturality of the pushforward construction, we have
  $\Psi(\PHk(f_t)) \cong (R(\phi_t^H)^{-1}_*) \Psi(\PHk(f))$.
  The sheaf functor $R(\phi_t^H)^{-1}_*$ is the Fourier--Mukai functor
  with kernel the graph sheaf $\mathcal{K}^H_t$ constructed by
  Theorem~\ref{thm:gks}~(ii), giving the intertwining
  equation~\eqref{eq:intertwining}.
\end{proof}

\begin{corollary}[Microlocal stability of $\PHk$ under modular flow]
  \label{cor:microlocal-stability}
  The persistence barcodes $\mathcal{B}(\PHk(f))$ are preserved, up to
  isomorphism and up to the convolution distance bound
  $\dconv \leq \|\phi_t^H - \id\|_{C^0}$, under the action of the
  modular Hamiltonian flow $\phi_t^H$.  In particular:
  \begin{enumerate}
    \item[(i)]
      The barcode decomposition of $\PHk$ is invariant under any
      Hamiltonian isotopy $\Phi$ that is non-displacing with respect to
      the microsupport of $\Psi(\PHk(f))$, by
      Theorem~\ref{thm:gks}~(iii).
    \item[(ii)]
      The ECI Euler-characteristic shift $\langle \chi_E(\nu)
      \rangle_{f_\mathrm{NL}} - \langle \chi_E(\nu) \rangle_0$ of axiom A6
      is stable under modular perturbations of the density field to
      first order in $\|\phi_t^H - \id\|$.
  \end{enumerate}
\end{corollary}

\begin{proof}
  Part~(i) is immediate from the microsupport non-displaceability
  criterion of Theorem~\ref{thm:gks}.  Part~(ii) follows from the
  stability bound in Theorem~\ref{thm:main}~(E3) applied to $f$ and
  $f_t = f \circ (\phi_t^H)^{-1}$, noting that
  $\|f_t - f\|_{L^\infty} \leq L(f) \cdot \|\phi_t^H - \id\|_{C^0}$
  for a Lipschitz constant $L(f)$ of $f$.
\end{proof}

% ==================================================================
\section{Discussion}
\label{sec:discussion}
% ==================================================================

\subsection{Closure of ECI axiom A6}

The main theorem (Theorem~\ref{thm:main}) and its corollary
(Corollary~\ref{cor:A6-derived}) establish that ECI's axiom A6 ---
the use of $\PHk$ on cosmological density super-level sets --- is a
consequence of microlocal sheaf theory under the hypotheses M1
(constructibility) and M2 (tameness).  This closes one of ECI's
standing axioms: $\PHk$ is no longer postulated but derived.

The derivation has the following concrete consequences for ECI:
\begin{enumerate}
  \item[(i)]
    The stability of $\PHk$ under $L^\infty$-perturbations of the
    density field $f$ (Theorem~\ref{thm:main}~(E3)) implies that the
    ECI $f_\mathrm{NL}$ estimator is robust to observational noise at
    the same rate as the classical bottleneck stability theorem
    \cite{CSGO2016}.
  \item[(ii)]
    The modular Hamiltonian intertwining
    (Theorem~\ref{thm:modular-sheaf-quantization}) implies that the
    ECI modular flow acts on the persistence sheaf via a precise
    Fourier--Mukai functor, giving a structural prediction: the barcodes
    of $\PHk$ transform equivariantly under modular evolution.
  \item[(iii)]
    The higher-dimensional generalisation (Remark~\ref{thm:main}) shows
    that the 3D cosmological $\PHk$ is controlled by the microsupport of
    $Rf_*\kk_X$ in $T^*\R^3$, providing a direct microlocal
    interpretation of cosmic web topology.
\end{enumerate}

\subsection{Relation to Berkouk--Ginot--Oudot and Berkouk--Petit}

The level-sets persistence theory of Berkouk--Ginot--Oudot
\cite{BGO2019} provides a further functorial equivalence: they prove
that \emph{level-sets} persistence modules (arising from Reeb graphs of
$f: X \to \R$) are equivalent, via a pseudo-isometric functor, to
derived pushforward sheaves $Rf_*\kk_X$ in $\Dbc(\Sh_\kk(\R))$.  This
extends our Theorem~\ref{thm:main} to the biparameter (birth and
death) setting, and connects the ECI density field observable to a 2D
persistence barcode stratification.

Berkouk--Petit \cite{BP2021} show that the observable category --- the
quotient of persistent modules by the ephemeral ones --- is equivalent to
the category of $\gamma$-sheaves, providing an independent categorical
path to the same equivalence.

\subsection{Relation to the Guillermou--Kashiwara--Schapira non-displaceability results}

The non-displaceability applications of Theorem~\ref{thm:gks} --- the
original motivation of \cite{GKS2012} --- have a cosmological
interpretation in ECI: a Hamiltonian isotopy that represents a modular
perturbation of the density field cannot ``displace'' the persistence
sheaf if the microsupport of the latter is non-displaceable.  This
constrains the topology-changing capacity of modular dynamics.

\subsection{Open problems}

\begin{enumerate}
  \item[(O1)]
    \textbf{Multiparameter persistence.}
    The present paper treats only 1-parameter persistence (sublevel sets
    of a single function $f$).  ECI's density field in 3+1 dimensions
    naturally leads to a 2-parameter persistence module
    $(f_1, f_2) = (\delta, \nu)$ (density and threshold).  The
    microlocal sheaf interpretation of 2-parameter persistence is an
    active area \cite{BGO2019}; extending our main theorem to this
    setting is the natural next step.

  \item[(O2)]
    \textbf{Constructibility hypotheses.}
    Hypothesis M1 (constructibility) is satisfied for algebraic and
    sub-analytic functions $f$, which covers the mathematical models of
    density fields.  Whether cosmological density fields from N-body
    simulations are ``tame'' in the required sense is a question of
    applied microlocal analysis.

  \item[(O3)]
    \textbf{Categorical upgrade.}
    The equivalence in Theorem~\ref{thm:main} is stated at the level of
    objects and morphisms in $\Dgamma(V;\kk)$.  A full $(\infty,1)$-categorical
    version, compatible with the higher-categorical structure of the
    crossed product algebra of ECI, would provide a stronger foundation
    for the ECI framework.

  \item[(O4)]
    \textbf{Spectral sequences and type-II algebras.}
    The modular Hamiltonian sheaf quantization
    (Theorem~\ref{thm:modular-sheaf-quantization}) suggests a sheaf-theoretic
    approach to the spectral sequence relating the Tomita--Takesaki modular
    theory of $\mathcal{M}$ to the derived category $\Dbc(\Sh_\kk(M))$.
    This is a deep problem connecting operator algebras and microlocal
    geometry.
\end{enumerate}

% ==================================================================
% Bibliography
% ==================================================================

\begin{thebibliography}{BGO19}

\bibitem[BG18]{BG2018}
N.~Berkouk and G.~Ginot,
\textit{A derived isometry theorem for constructible sheaves on $\mathbb{R}$},
arXiv:1805.09694 [math.AT], 2018.
% [V] verified 2026-05-04 via arXiv abstract

\bibitem[BGO19]{BGO2019}
N.~Berkouk, G.~Ginot, and S.~Oudot,
\textit{Level-sets persistence and sheaf theory},
arXiv:1907.09759 [math.AT], 2019.
% [V] verified 2026-05-04 via arXiv abstract

\bibitem[BP21]{BP2021}
N.~Berkouk and F.~Petit,
\textit{Ephemeral persistence modules and distance comparison},
Algebr. Geom. Topol. \textbf{21} (2021), 247--277,
arXiv:1902.09933 [math.AT].
% [V] verified 2026-05-04 via arXiv abstract; journal ref confirmed

\bibitem[BW76]{BisognanoWichmann1976}
J.~J.~Bisognano and E.~H.~Wichmann,
\textit{On the duality condition for a Hermitian scalar field},
J. Math. Phys. \textbf{16} (1975), 985--1007;
\textit{On the duality condition for quantum fields},
J. Math. Phys. \textbf{17} (1976), 303--321.

\bibitem[CB15]{CrawleyBoevey2015}
W.~Crawley-Boevey,
\textit{Decomposition of pointwise finite-dimensional persistence modules},
J. Algebra Appl. \textbf{14} (2015), 1550066,
arXiv:1210.0819 [math.RA].

\bibitem[CLPW23]{CLPW2023}
V.~Chandrasekaran, R.~Longo, G.~Penington, and E.~Witten,
\textit{An algebra of observables for de~Sitter space},
J. High Energy Phys. \textbf{2023}:082,
arXiv:2206.10780 [hep-th].

\bibitem[CSGO16]{CSGO2016}
F.~Chazal, V.~de~Silva, M.~Glisse, and S.~Oudot,
\textit{The Structure and Stability of Persistence Modules},
SpringerBriefs in Mathematics, 2016,
arXiv:1207.3674 [cs.CG].

\bibitem[Cur14]{Curry2014}
J.~Curry,
\textit{Sheaves, Cosheaves and Applications},
Ph.D. thesis, University of Pennsylvania, 2014,
arXiv:1303.3255 [math.AT].

\bibitem[DEHK25]{DEHK2025}
J.~De Vuyst, S.~Eccles, P.~A.~H\"ohn, and J.~Kabel,
\textit{Crossed products, extended phase spaces and the resolution of
Entanglement Island paradoxes},
arXiv:2412.15502 [hep-th], 2024.

\bibitem[ECI26]{ECI2026}
K.~Remondi\`ere,
\textit{ECI --- Entanglement, Complexity, Information: a framework paper
threading six programs into one predictive canvas},
v6.0.49, Zenodo DOI \href{https://doi.org/10.5281/zenodo.19686398}{10.5281/zenodo.19686398},
2026.

\bibitem[EH10]{EdelsbrunnerHarer2010}
H.~Edelsbrunner and J.~L.~Harer,
\textit{Computational Topology: An Introduction},
American Mathematical Society, 2010.

\bibitem[GKS12]{GKS2012}
S.~Guillermou, M.~Kashiwara, and P.~Schapira,
\textit{Sheaf quantization of Hamiltonian isotopies and applications to
non-displaceability problems},
Duke Math. J. \textbf{161} (2012), 201--245,
arXiv:1005.1517 [math.SG].
% [V] verified 2026-05-04 via arXiv abstract and journal publication data

\bibitem[KS90]{KS1990}
M.~Kashiwara and P.~Schapira,
\textit{Sheaves on Manifolds},
Grundlehren der Mathematischen Wissenschaften, vol.~292,
Springer-Verlag, Berlin, 1990.

\bibitem[KS18]{KS2018}
M.~Kashiwara and P.~Schapira,
\textit{Persistent homology and microlocal sheaf theory},
J. Appl. Comput. Topology (2018, accepted),
arXiv:1705.00955 [math.AT].
% [V] verified 2026-05-04 via arXiv abstract; journal acceptance confirmed

\end{thebibliography}

\end{document}
